% 04_scheduler
\section{调度与大小核}

调度器决定每个可运行线程在何时、在哪个核上运行；在多核系统上，一个线程被放到哪条运行队列，就决定它在哪个核上执行。StarryOS 的 axtask 此前不做跨核放置：新建的线程固定留在派生它的核，被唤醒的线程回到唤醒者所在的核，运行期也没有负载均衡。RK3588 由 4 个 A55 小核与 4 个 A76 大核组成，整机八核的算力因此始终集中在启动所在的单个小核上，sysbench cpu 无论开多少线程都停在 \nLadderBase{} ev/s 附近，八线程落后同机 Linux 约三十倍。现已在 axtask 上补齐跨核派生分发、占用感知的派生与唤醒放置、Linux-CFS 式的 \texttt{wake\_affine} 就近唤醒，并为 \texttt{CLONE\_VM} 线程负载补上一条无锁的用户拷贝快路径；八线程 sysbench cpu 升到 \keyres{\nCpuTEight{} ev/s，为 Linux \nCpuTEightLin{} 的 \nCpuTEightRatio{} 倍}，单核算力与 Linux 对等，核间驻留均匀。本节先说明调度器与运行队列的放置机制、以及大小核给放置带来的约束，再逐一说明四层新增机制的原理与实现，最后给出对等阶梯、核驻留与单核算力的板级结果。

\figphl{../figures/out/Fsched.png}{axtask 放置管线，青色为新增实现；顶部为待放置线程，各核持私有运行队列。派生分发在亲和性允许的核间轮转，单原子 \texttt{occ} 占用计数供派生与唤醒共用并每 tick 自愈，\texttt{wake\_affine} 在唤醒者空闲时把被唤醒者就近交回，共同把负载在大小核间铺开}{fig:sched-arch}

\subsection{调度器与运行队列}

调度器的职责是把有限的 CPU 时间分配给众多可运行线程。axtask 为每个 CPU 维护一条独立的运行队列：核空闲时从本核队列按调度策略取下一个线程执行，线程时间片耗尽或主动阻塞时再重新入队。运行队列为每核私有，一个线程当前排在哪条队列上，就决定它在哪个核上运行。把线程指派到某条运行队列的这一步称为放置。

放置发生在两个时刻，各对应 \texttt{run\_queue.rs} 里的一处决策。其一是派生放置：\texttt{fork}／\texttt{clone} 新建一个线程时，须为它选一条初始运行队列（\texttt{select\_run\_queue}）。其二是唤醒放置：一个阻塞在某事件上的线程被唤醒、重新变为可运行时，须为它选一条运行队列重新入队（\texttt{select\_wake\_run\_queue}）。这两处决策共同决定负载在各核之间如何分布：派生放置铺开新建的工作线程，唤醒放置决定周期性阻塞与唤醒的线程醒来后落在哪。放置一旦失衡，即便有空闲核也无线程可取，整机吞吐会退化到被占满的少数核的算力水平。

\subsection{大小核与既有 axtask 的放置}

RK3588 采用大小核结构，8 个核分属两簇：cpu0 至 cpu3 为 A55 小核，cpu4 至 cpu7 为 A76 大核，A76 的每周期吞吐与频率上限均高于 A55。启动核为 cpu0，属 A55 小核簇。异构给放置多加一层约束：同一线程放在 A76 与放在 A55 上吞吐不同，只有把可并行的线程铺满全部八核、并让繁重线程优先落在大核，才能发挥整机算力。

既有 axtask 在两处放置点都不跨核迁移，运行期也不做负载均衡，三点叠加使整机负载始终集中在单个核上。

\paragraph{派生不跨核} 派生放置把新线程固定在执行 \texttt{clone} 的那个核上。\texttt{fork}／\texttt{clone} 几乎总在启动核 cpu0 上发起，新建线程因而全部堆在 cpu0 这一个 A55 小核。

\paragraph{唤醒回汇聚} 唤醒放置把被唤醒者交回唤醒者所在的核。栅栏同步的一组工作线程由同一个释放者唤醒，醒来后重新汇聚回释放者的核，无法分散。

\paragraph{无负载均衡} 运行期没有任何均衡逻辑。一条队列排满、其余七条全空时，没有机制把线程迁到空闲核。

三点叠加的结果是整机只在单个小核上排队，smp8 与 smp4 的读数逐字节相同，sysbench cpu 无论线程数多少都停在 \nLadderBase{} ev/s 附近，八线程落后同机 Linux 约三十倍。这一差距纯由放置造成：一组强制每核各绑定一个任务的对照读到接近八核理论上限的吞吐，确认单核算力本身并无问题，缺的是把线程铺开的放置逻辑。现已在这条放置路径上新增四层机制，自派生到唤醒依次生效，其结构与落点见图~\ref{fig:sched-arch}。

\subsection{跨核派生与唤醒分发（\#1656）}

新增的第一层是跨核派生与唤醒分发，提交为 \#1656。派生时不再固定到启动核，而\ours{在亲和性允许的核之间轮转分发}；唤醒时优先选线程上一次运行的核，保住缓存热度，不再一律交回唤醒者。轮转须只落在已上线的运行队列上，为此新增一个 \texttt{RUN\_QUEUE\_ONLINE} 位图门控，避免早期启动阶段把线程分派到尚未初始化的次核队列而触发 \texttt{DataAbort}。派生分发与唤醒改选两半缺一不可：仅做派生轮转时，栅栏唤醒的工作线程仍会被重新汇聚回唤醒者。

轮转把真并发暴露出来后，几处此前被启动核串行化所掩盖的缺陷随之显现，均按 Linux 语义修正并随 \#1656 一并验证。一处是 SMP 下的 ptrace 缺陷：\texttt{wait\_ptrace\_resume} 会因一次善意的跨核 \texttt{interrupt()} 放弃 trace-stop，使刚执行 \texttt{TRACEME} 的子进程越过 \texttt{raise(SIGSTOP)} 继续运行。一处是 riscv64 的远程 kick 合并竞态：\texttt{REMOTE\_RESCHEDULE\_PENDING} 会压制掉排在在途 reschedule 之后落核任务的 IPI，把它搁在无 tick 的 idle 上，以 \texttt{force\_kick\_remote\_cpu} 补发 IPI 修复。一处是 \texttt{pivot\_root} 后关机时的卸载 panic：init 的 \texttt{unmount\_all} 在共享挂载命名空间未静默时 panic，改为尽力卸载，并增设一把对应 Linux \texttt{namespace\_sem} 的全局 \texttt{MOUNT\_OP\_LOCK}。此外补上一处 POSIX 亲和性缺口：\texttt{do\_clone} 从不复制父线程的 cpumask，使一个 \texttt{taskset -c 0} 限定的 sysbench 仍可能运行到 A76 上，修正为在克隆时令子线程继承父线程的亲和性掩码。

跨核分发把上板 sysbench cpu \keyres{从平线 \nLadderBase{} 抬到八线程 \nLadderA{} ev/s}，确认此前的差距纯粹来自放置。

\subsection{占用感知的放置}

派生轮转把线程散开，但在异构核上分布仍不均匀，因为轮转看的是队列长度 \texttt{nr\_running}，而它不计入当前正在运行的任务。一个 A76 正运行一个计算线程、队列却为空时，\texttt{nr\_running} 读作 0，这个繁忙的大核被当成空闲，连续派生的兄弟线程于是全部被放置到该核，形成聚集。

准确的占用需要把正在运行的任务也计入。一个直接的做法是在 \texttt{nr\_running} 之外再加一个 \texttt{RQ\_RUNNING} 标志、两者相加得占用；但这两个量无法被远端放置者一致地读到，属 seqlock 类竞态：运行态转空闲的切换从不触碰 \texttt{nr\_running}，那次置 false 的存储没有配对的 Release，一次带序存储的尝试因此反而回归并被回退。最终以\ours{单个每核 \texttt{occ} 原子计数}解决：\texttt{occ} 仅在线程进入和离开运行态时（在 \texttt{switch\_to} 内）增减，把正在运行的任务计入负载，远端放置者只读一个原子量即可得到一致的占用。

发行版构建会移除针对 \texttt{occ} 下溢的 \texttt{debug\_assert}，使 \texttt{occ} 在长时间运行后缓慢上漂，大核被误判为繁忙。为此在调度 tick 上加一层自愈：\texttt{scheduler\_timer\_tick} 每约 10 ms 重算一次占用：取 \texttt{nr\_running} 与正在运行的非 idle 任务数之和，并配一个饱和递减，把漂移校正回真实占用。

唤醒放置同样接入 \texttt{occ}，对齐 Linux 的 \texttt{select\_idle\_sibling}：仅当目标核空闲（\texttt{occ==0}）时保留线程上一次运行的核，否则改选最空闲的核（\texttt{select\_least\_loaded}）。此前唤醒把被唤醒者固定回 \texttt{last\_cpu} 或唤醒者，栅栏释放的一阵唤醒会把八个线程聚到三个大核上，接入占用信号后八个线程才在全部八核上摊开。占用感知的派生与唤醒放置在大核频率提升带来的 \nLadderB{} ev/s 基础上\keyres{把八线程 sysbench cpu 推进到 \nCpuTEight{} ev/s}，并使核间驻留由集中转为均匀。

\subsection{完整工作窃取均衡器的反面结论}

在 \#1656 基础上另前向移植过一套完整的容量感知工作窃取均衡器，含空闲核 idle-pull、繁忙核 push 与唤醒重定向，并从设备树读取每核容量（\texttt{capacity-dmips-mhz}，A76 为 1024、A55 为 530）按容量折算负载。主机测试全绿，但在板上它使目标负载回归：单线程吞吐由 350 压到 161，饱和时并无收益。其机制为迁移抖动，七个空闲核把唯一的工作线程在核间反复迁移，无一核心得以保持缓存热度。配套移植的一次 Linux 式 \texttt{nr\_balance\_failed} 升级会让空闲 A55 抢走热 A76 上的任务，同样回退。运行期迁移（\texttt{idle\_pull\_once}／\texttt{try\_push\_balance}）因此仅以特性门控保留、默认关闭，出厂默认只采用容量感知的派生与唤醒放置这一安全子集。

\subsection{wake\_affine 就近唤醒}

清除分页相关的 EFAULT（详见内存一节）后，整机 schbench 暴露出约 1 ms 的跨核唤醒下限：m1t4 的 wakeup p50 约 \nWakeBefore{} µs，同机 Linux 仅 \nWakeLin{} µs。

一次跨核唤醒的开销主要在处理器间中断投递。线程 A 在核 X 上唤醒线程 B、而 B 被放到另一个核 Y 时，核 X 须向核 Y 发一次 reschedule 的 GIC SGI（一种处理器间中断）通知它有新线程可运行，核 Y 收到后才切入 B。\texttt{wake\_affine} 针对生产者唤醒单个消费者这类同步唤醒：当唤醒者近乎空闲（\texttt{occ<=1}）时，\ours{把被唤醒者交回唤醒者所在的核，就近运行}，省去这次跨核 SGI，且共享缓存仍热。这一层对齐 Linux CFS 的 \texttt{wake\_affine}，默认启用；它把 m1t4 的 wakeup p50 \keyres{由约 \nWakeBefore{} µs 降到 \nWake{} µs，接近 Linux 的 \nWakeLin{} µs}，判定持平。至于栅栏同步的一组并发工作线程被同一释放者唤醒的情形，就近交回并不适用，留待下文的唤醒铺开处理。

这约 1 ms 全部来自 IPI 投递本身。带内分类计时显示本地唤醒 p50 为 2 µs，跨核项约 1048 µs，而中断处理到切入仅 2 µs，且 87\% 的 reschedule SGI 指向真正处于 WFI 的核心。RK3588 的 reschedule GIC SGI 无法及时唤醒处于 WFI 的空闲核，该核要等自身约 1 到 2 ms 的 oneshot 定时器到期才前进。\texttt{wake\_affine} 以就近唤醒绕开这条慢投递。对确需跨核唤醒的情形，idle-poll（WFI 前自旋约 50 µs）把跨核唤醒中位数压低约 250 倍，下文的唤醒铺开正依赖这一点。该下限属固件与硬件层面的外因，软件仅能缓解，无从消除。

\subsection{并发工作线程的唤醒铺开}

\texttt{wake\_affine} 的就近唤醒适合生产者唤醒单个消费者，对栅栏释放的一组并发工作线程却适得其反：它们由同一个释放者唤醒，就近交回会把整组重新聚拢到释放者的核，抵消派生放置本已把它们铺开的效果。这类聚拢在两线程 mem\_bw2 与 schbench 上显现，一组本应各占一个大核的工作线程挤在同一个核上。直接套用 Linux 的 \texttt{wake\_affine\_idle} 加 \texttt{select\_idle\_sibling} 把每次唤醒都改为铺开，在缺 idle-poll 时反而更慢，因每次铺开都要支付前述约 1 ms 的跨核 SGI 代价；以固定时长自旋的 idle-poll 抹平这笔代价之后，铺开转为净收益。出厂放置配置据此\ours{一并启用 idle-poll，并以按线程 CPU 占用的 EWMA 区分两类唤醒}：判为 CPU 密集的并发工作线程无条件经 \texttt{select\_idle\_sibling} 铺开到空闲兄弟核，生产者消费者仍走 \texttt{wake\_affine} 就近；两者都在唤醒时决定落核，仍属唤醒放置的安全子集，不涉及前述已回退的运行期迁移。板级 A/B 中，schbench m1t4 的 RPS 由约 \nSchRpsBefore{} 抬到 \nSchRps{}，为同机 Linux 约 \nSchRpsLin{} 的 \nSchRpsRatio{} 倍；两线程 mem\_bw2（2 个 A76）由 0.34 抬到 \nMemBwTwo{} 倍 Linux，栅栏铺开由单核扩到全八核；sysbench mutex 八线程的墙钟由 \nMutexBefore{} s 降到约 \nMutexAfter{} s，对 Linux 的比值为 \nMutexRatio{} 倍、残差在锁路径本身（详见后文残余差距）。上述数值取自 2026-08-12 板级消融的三次重复中位数。

\subsection{\texttt{-T} 线程负载的用户拷贝快路径}

hackbench 的 \texttt{-T}（\texttt{CLONE\_VM} 线程）此前比 \texttt{-P}（fork 进程）慢 3.6 到 4.2 倍，方向与 Linux 相反。根因是一把共享、独占且会睡眠的 \texttt{Arc<Mutex<AddrSpace>>}：每一次用户内存拷贝前的地址校验都要锁住整个地址空间，且这把锁的持有跨越一次页表走查。\texttt{CLONE\_VM} 线程共享同一个地址空间，全部在这把锁上串行；fork 出的进程各自持有私有地址空间，互不相干，故 \texttt{-T} 慢而 \texttt{-P} 快。

新增的快路径对齐 Linux 的 \texttt{access\_ok} 模型，在取锁之前先做一次\ours{无锁的硬件页表探测}。\texttt{access\_ok} 的职责是在内核拷贝用户内存前确认目标地址在当前进程地址空间内可读或可写，既有实现靠锁住地址空间并软件走查页表来判定。aarch64 提供地址翻译指令 \texttt{AT}：\texttt{AT S1E0R} 让硬件按 stage-1、EL0、读权限对一个虚拟地址做一次翻译与权限检查，把结果写入 \texttt{PAR\_EL1}，其 F 位（bit 0）为 0 表示该地址在 EL0 可读、为 1 表示不可达；\texttt{AT S1E0W} 同理判可写。发出 \texttt{AT} 后以 \texttt{isb} 等翻译完成，再读 \texttt{PAR\_EL1}。\texttt{components/axcpu} 据此实现 \texttt{user\_access\_ok\_page}：关中断下对每个待访问页发一次 \texttt{AT} 探测（上限 \texttt{FASTPATH\_MAX\_PAGES=16} 页），全部可达则跳过取锁直接拷贝，任一不可达再回落到取锁的慢路径，接入 \texttt{check\_region} 与 \texttt{prepare\_user\_memory}。前提均经核实：PAN 关闭（\texttt{SCTLR\_EL1.SPAN=1}）使 EL1 得以按 EL0 权限访问用户地址，TBI0 关闭，COW 破坏集与 \texttt{AT} 接受集证明不相交。一次对抗审查发现一个校验绕过：一个会回绕的 \texttt{UserPtr} 在 \texttt{page\_start>page\_end} 时跳过探测循环直接返回可达，特性开启时构成校验绕过；以 \texttt{checked\_add} 修正后复核确认无绕过。

板级 A/B（同一 HEAD，仅特性开关不同）中，hackbench \keyres{\texttt{-T} g10 提升 \nHackT{} 倍}，\texttt{-T} 对 \texttt{-P} 的比值由 3.6 到 4.2 倍收敛到约 1，恢复 Linux 的形态，\texttt{-P} 自身也加快 1.5 到 2 倍，已并入出厂配置。另有一处 COW 引用计数溢出导致的 fork EFAULT 缺陷，其修复详见内存一节。

\subsection{对等阶梯与核驻留}

固定频率隔离出放置各层的净贡献，构成 sysbench cpu 的对等阶梯：单核地板（放置关闭）约 \nLadderFloor{} ev/s，补齐跨核分发即跃至约 \nLadderPlace{} ev/s（为 Linux 的 \nCpuTEightRatio{} 倍，放置是决定性的一层）；占用感知放置与出厂配置（含唤醒铺开）在此量级上持平，分别约 \nLadderOcc{} 与 \nCpuTEight{} ev/s，为 Linux \nCpuTEightLin{} 的 \nCpuTEightRatio{} 倍，判定持平。\texttt{wake\_affine} 单独的就近唤醒会把八线程回落到约 \nLadderCoalesce{} ev/s（为压低唤醒延迟所必需），唤醒铺开随即将其恢复。

\figph{../figures/out/F04.png}{sysbench cpu 对等阶梯（频率固定）。单核地板（放置关闭）约 \nLadderFloor{} ev/s，补齐跨核分发跃至约 \nLadderPlace{} ev/s，占用感知放置与出厂配置在此量级持平，分别约 \nLadderOcc{} 与 \nCpuTEight{} ev/s，为 Linux \nCpuTEightLin{} 的 \nCpuTEightRatio{} 倍}

放置均衡前，八线程会塌陷到少数大核。此前 \texttt{/proc/stat} 的每核驻留是聚合值除以核数的虚构读数，无法显出失衡；将其\ours{改为真实的每核忙碌 tick}（经 \texttt{ax\_task::cpu\_busy\_ticks} 暴露）后，一次八线程测量给出 \texttt{[cpu0 83, cpu1 0, cpu2 0, cpu3 0, cpu4 40, cpu5 2071, cpu6 1560, cpu7 851]}，A55 集群近乎空闲，八个线程集中在三个 A76 上，据此确认症结在聚集，空闲的 A55 本可分担却无线程落入。唤醒接入占用信号后，同一测量读到均匀驻留 \texttt{[951, 988, 959, 968, 944, 954, 929, 891]}。

\figph{../figures/out/F05.png}{核驻留直方图。放置均衡前八个线程集中在三个 A76 核、A55 集群近乎空闲，唤醒接入占用信号后在 A55 与 A76 八核间均匀铺开}

\subsection{单核算力与残余差距}

把聚合读数拆分到单核，算力已然对齐。A55 单核为 \nAfive{} ev/s，Linux 为 \nAfiveLin{}，比值 \nAfiveRatio{} 倍，小幅领先；A76 单核为 \nAseven{} ev/s，Linux 为 \nAsevenLin{}，比值 \nAsevenRatio{} 倍，持平；综合每核\keyres{约为 Linux 的 \nPerCoreRatio{} 倍}。八线程聚合 \nCpuTEight{} ev/s 为 Linux \nCpuTEightLin{} 的 \nCpuTEightRatio{} 倍，判定持平。

\begin{table}[H]\centering\small
\begin{tabular}{@{}p{3.6cm}rrrl@{}}
\toprule
口径 & StarryOS & Linux & 比值 & 判定 \\
\midrule
单核 A55（ev/s） & \nAfive{} & \nAfiveLin{} & \nAfiveRatio{}× & 小幅领先 \\
单核 A76（ev/s） & \nAseven{} & \nAsevenLin{} & \nAsevenRatio{}× & 持平 \\
cpu 8 线程（ev/s） & \nCpuTEight{} & \nCpuTEightLin{} & \nCpuTEightRatio{}× & 持平 \\
\bottomrule
\end{tabular}
\caption{每核与八线程 sysbench cpu 对等，p50 取 N$\geq$5}
\end{table}

其距满对等的残余来自频率上限：大核 DVFS 提升到 \nDvfsAseven{} MHz（上探 \nDvfsAsevenTop{} MHz）、小核 \nDvfsAfive{} MHz 后，仍受 BL31 安全 CRU 与温控缺失所限，为外因所致，详见频率一节。锁竞争是后续继续优化的方向：放置与频率均已对齐，mutex 与 futex 的残差集中在锁自身的路径。总体而言，调度在结构与单核算力上与 Linux 对等，放置、跨核唤醒延迟与 \texttt{-T} 扩展均已恢复到 Linux 的形态。
